\chapter{Attention、KV Cache 与内存流量}

Attention 是 Transformer 的核心。推理时，它不仅包含矩阵乘，还包含 Softmax、mask、KV Cache 读写和随序列长度增长的内存访问。理解 Attention 性能，必须区分 prefill 和 decode：prefill 处理整段输入，矩阵形状较大；decode 每次生成一个 token，batch 和查询长度通常很小，但要读取历史 KV Cache。

\section{Attention 的计算图}

单头 Attention 可以写成 $O=\mathrm{softmax}(QK^T / \sqrt{d})V$。多头 Attention 只是把 hidden dimension 切成多个 head 并并行计算。这个公式里有两个主要矩阵乘：$QK^T$ 产生注意力分数，Softmax 归一化后再乘 $V$ 得到输出。若序列长度为 $S$，head dimension 为 $D$，那么分数矩阵大小是 $S \times S$，长序列时中间矩阵非常大。

prefill 阶段，$Q$、$K$、$V$ 都来自当前输入序列，计算类似批量矩阵乘。decode 阶段，新 token 只有一个或少量 query，但必须与历史所有 key 做点积，再对所有历史位置 Softmax，然后加权所有 value。此时计算量随历史长度线性增长，KV Cache 的读取成为关键。

\section{KV Cache 的作用}

如果每生成一个 token 都重新计算所有历史 token 的 $K$ 和 $V$，成本会非常高。KV Cache 把历史 token 的 key 和 value 保存下来，新 token 到来时只计算它自己的 $Q,K,V$，然后用新 $Q$ 与缓存中的所有历史 $K$ 做 Attention。这样减少了重复计算，但增加了持续增长的内存占用和读取流量。

KV Cache 的布局直接影响性能。常见维度包括 layer、batch、head、sequence、head dimension。哪一维连续，决定 decode 时读取一个 head 的历史 key 是否顺序，多个 head 是否容易并行，追加新 token 是否连续写入。布局还影响量化 KV Cache、分页 KV Cache 和多请求 batching。

\begin{lstlisting}[language=C++,caption={教学版 KV Cache 索引：明确每个维度的地址贡献}]
#include <cstddef>
#include <cstdint>

struct KvShape {
    std::int32_t layers;
    std::int32_t batch;
    std::int32_t heads;
    std::int32_t max_seq;
    std::int32_t head_dim;
};

std::size_t kv_index(const KvShape& shape,
                     std::int32_t layer,
                     std::int32_t batch_id,
                     std::int32_t head,
                     std::int32_t token,
                     std::int32_t dim) {
    std::size_t index = static_cast<std::size_t>(layer);
    index = index * shape.batch + static_cast<std::size_t>(batch_id);
    index = index * shape.heads + static_cast<std::size_t>(head);
    index = index * shape.max_seq + static_cast<std::size_t>(token);
    index = index * shape.head_dim + static_cast<std::size_t>(dim);
    return index;
}
\end{lstlisting}

这个布局让同一个 token 的 head dimension 连续。如果 decode 时要对固定 head 扫描所有历史 token，并且每个 token 读取完整 head dimension，那么 token 维和 dim 维的组合访问是规则的。但如果 batching、分页或量化改变布局，就要重新分析地址序列。

\section{从完整矩阵到在线 Softmax}

朴素 Attention 会先算完整分数矩阵，再做 Softmax，再乘 $V$。这对长序列很昂贵，因为分数矩阵写入和读出都很大。在线 Softmax 的思想是按块处理分数，维护当前最大值和归一化分母。当新块到来且最大值变化时，旧的累加结果按比例缩放。这样可以不把完整分数矩阵写回内存。

这类算法说明，Attention 优化不是简单“调用更快 GEMM”。它要把矩阵乘、Softmax 和乘 $V$ 融合成数据流问题。分块大小要适配缓存，统计量要保持数值稳定，mask 和 causal 约束要正确处理，尾部和不同序列长度也要覆盖。

\section{decode 阶段的瓶颈}

decode 阶段常见瓶颈是内存带宽和小矩阵效率。每生成一个 token，都要读取每层的 KV Cache。随着上下文变长，读取历史 KV 的字节数增长，而 query 很少，矩阵乘形状偏瘦，难以达到大 GEMM 的峰值效率。此时优化方向包括 KV Cache 量化、分组查询注意力、分页管理、连续 batching、减少中间写回和更好的线程切分。

\begin{warningbox}
不要用 prefill 的吞吐推断 decode 的速度。prefill 更像大矩阵计算，decode 更像长历史缓存扫描加小矩阵计算。两者的瓶颈、批处理策略和优化手段都不同。
\end{warningbox}

\section{源码阅读入口}

阅读 llama.cpp、vLLM、oneDNN Graph 或 XNNPACK 这类项目时，Attention 路径可以按四个问题定位。第一，$Q,K,V$ 投影使用什么矩阵乘路径。第二，KV Cache 的物理布局是什么，追加和读取如何索引。第三，Softmax 是否与分数计算或乘 $V$ 融合。第四，prefill 和 decode 是否走不同 kernel。只要这四个问题清楚，复杂代码就能分层理解。

\begin{exercisebox}
给定 batch 为 1、head 数为 32、head dimension 为 128、上下文长度为 4096 的 decode 场景，估算一层读取 K 和 V 至少需要多少字节。再乘以层数，观察为什么 KV Cache 带宽会成为长上下文推理的重要问题。
\end{exercisebox}
